\begin{afterword}
\summaryhead

The post looks back on early quantum physics for a lesson. Its nominee for the most basic error is that the founders forgot they were themselves made of particles. They did not see that a detector, or their own knowledge of an electron's path, is just more particles in different places, so they brought consciousness into the theory. Knowing things feels simpler than brain states, as Thor once felt simpler than Maxwell's equations. A theory in which observation plays a part could not be written as a program, because nobody said in physical terms when an observation occurs.

The needed insight looks philosophical, but physicists, not philosophers, cleared up the mystery. In general, the post says, science and not the consensus of philosophers tells us which philosopher was right. ``Philosophy'' names thinking not yet taught or reduced to calculation. It is done well only within a science, and philosophers should learn the science and expect experiment to settle their arguments.

\responsehead

The post draws a lesson about philosophy from one history, told without naming any of the physicists it faults. The views that gave consciousness a role in quantum measurement are usually traced to von Neumann, London and Bauer, and Wigner. Von Neumann did not leave the instruments out of the physics. He allowed that the measuring instrument could be described by ordinary quantum evolution, and he ended the regress the post itself describes with an ``abstract ego.'' Bohr, the main founder of the Copenhagen view, rejected any role for the observer's mind. The suppression of interference was studied from the start. So the charge that the founders ``forgot that they themselves were made out of particles'' does not fit von Neumann, who applied the theory to the instruments and still posited a cut.

The claim that physicists, not philosophers, found the answer fits its origin. Everett was a physics graduate student, and decoherence theory came from Zeh and Zurek. But ``solved the problem and cleared up the mystery'' conflicts with what the author wrote a month later, when ``Many Worlds, One Best Guess'' said there is ``no known reason for the Born probabilities.'' Philosophers of physics, among them Wallace, Greaves and Saunders, are part of the later work on the open questions. Whether many-worlds is right is still disputed among physicists.

The general rule, that science ``usually'' tells us which philosopher was right, does not fit the post's own case. In ``The Dilemma: Science or Bayes?'' the author rests many-worlds on simplicity, says the evidence is not in dispute, and says the physicists who reject it are following the rules of Science. On that account no experiment has decided the question.

The post has support it does not cite. The physicist John Bell made the same complaint about the undefined ``measurement'' in 1990, and Russell gave a version of the saying about sciences leaving philosophy in 1912.

In short: a lesson about philosophy drawn from one history whose physicists go unnamed, whose diagnosis does not fit von Neumann's treatment of measurement, and whose main case, by the author's account a month later, depends on simplicity rather than on the ``experimental results'' the post says tell the world which philosopher won.
\end{afterword}
